\honest{The Crackpot Offer}{Eliezer Yudkowsky}{2007}

When I was thirteen or maybe fourteen, I thought I had disproved Cantor's diagonal argument, the theorem that shows the real numbers outnumber the rational numbers. I dreamed of fame. \nb{The argument shows that no list indexed by the whole numbers contains every infinite sequence of 0s and 1s, or every set of whole numbers. The reals outnumber the rationals because the rationals can be listed that way and the reals cannot.}

My idea: write a whole number in binary, and the positions of its 1s give a set of whole numbers. The number 13 is 1101, which gives \{0, 2, 3\}. So the whole numbers could be matched with all the sets of whole numbers. A week later I applied the diagonal argument to my matching, and it found a set I had missed: the binary number $\ldots1111$. \nb{The mathematics is right. The matching covers exactly the finite sets. The diagonal argument, applied to it, gives the set of all whole numbers, since bit $n$ of $n$ is always 0.}

I was disappointed, and my first thought was that I would disprove the theorem someday. Then I saw that, having found my mistake, I had no more reason to doubt this theorem than any other. I was being offered the chance to become a math crank writing ``angry letters in green ink.'' I laughed and let it go. It was an unfair test for a child, I say, and I was lucky that I found the mistake myself and that it was simple.

The smarter you are, I suggest, the younger you may be when you first have what looks like a revolutionary idea. I wonder how many cranks were thirteen when they took their first wrong step.

It was a mistake, and that was all: not ``half right or even the tiniest fraction right.'' \nb{For a claimed disproof this is exact: a proof with a gap proves nothing. Many mistaken beliefs, such as a forecast or a plan, can be partly right, and the post gives no way to tell the two kinds apart.} Had I read it as partly right, I would have kept looking for a flaw, and might have found one.

My rule: ``Whenever you are tempted to hold on to a thought you would never have thought if you had been wiser, you are being offered the opportunity to become a crackpot.'' It is the clinging that makes a crackpot, even if you never tell anyone. \nb{A thought can be recognized as one ``you would never have thought if you had been wiser'' only after the mistake is found. The rule says nothing about how to tell, before then, a mistaken thought from a correct unpopular one.} Not every cloud has a silver lining. Say ``oops,'' and get on with your life.
